% Secao 10 - Limitacoes. Honesta, sem hedging redundante.

\section{Limitações}
\label{sec:limitacoes}

Quatro limites delimitam o que o relatório pode afirmar. O primeiro é a cobertura
da fonte: o SIH registra apenas a internação financiada pelo SUS, de modo que o
fluxo é cego à saúde suplementar. Em municípios de alta penetração privada, o
burden de fluxo superestima a dependência efetiva de hospitais distantes; é por
isso que a cobertura ANS entra como moderador na Seção~\ref{sec:divergencia}, e
por isso o achado central se concentra justamente onde a cobertura privada é
baixa e o SIH é mais representativo.

O segundo é a natureza endógena do fluxo. Diferentemente da distância, que é
geográfica e exógena ao comportamento, o burden de fluxo reflete decisões de
encaminhamento, oferta instalada e a própria capacidade de viajar. Isso é uma
virtude para mensuração --- o fluxo capta acesso revelado, não acesso potencial
--- mas impede qualquer leitura causal: F1 é uma medida de exposição, não um
tratamento. O relatório é, por construção, descritivo.

O terceiro diz respeito às aproximações geográficas. A distância usa centroides
municipais e a malha rodoviária, o que subestima deslocamentos intramunicipais
em capitais e não representa a viagem fluvial na Amazônia --- a fragilidade que a
própria comparação com o fluxo expõe. A estrutura etária usada na padronização da
Seção~\ref{sec:mortalidade} vem do Censo 2022 e é assumida constante na janela.

O quarto é a sensibilidade a escolhas de corte. A taxonomia de quadrantes usa o
quartil superior de cada medida como limiar de deserto; limiares absolutos
alternativos e as operacionalizações de fluxo F2 (isolamento em rede) e F3
(concentração de destinos) são reportados no apêndice, e concordam com o padrão
central. Nenhuma conclusão do corpo depende de um único limiar ou de uma única
definição de fluxo.
